%!TEX root = ../thesis.tex
%*******************************************************************************
%****************************** Third Chapter **********************************
%*******************************************************************************

\chapter{Methodology}\label{chap:methodology}

% **************************** Define Graphics Path **************************
\ifpdf
    \graphicspath{{Chapter3/Figs/Raster/}{Chapter3/Figs/PDF/}{Chapter3/Figs/}}
\else
    \graphicspath{{Chapter3/Figs/Vector/}{Chapter3/Figs/}}
\fi

\section{Methodology Overview}
This chapter outlines the research design used to build and evaluate the conversational assistant. The project followed a three-month iterative cycle (November 2024 -- January 2025) that combined requirements elicitation, data curation, prototype implementation, and empirical evaluation. The assistant targets psychiatric nurses and supports natural language queries about patient profiles. The overarching methodological goal was to determine whether a retrieval-augmented conversational workflow could deliver timely, accurate information while aligning with clinical expectations.

\section{Data Sources and Collection Methods}
Requirements gathering relied on qualitative interviews with three psychiatric nurses who each participated in a two-hour remote session. The interviews surfaced the daily responsibilities of nurses, critical pain points, and essential data elements required for safe care. Key needs included continuous patient observation, medication administration, collaboration with multidisciplinary teams, and rapid access to the handwritten ``Patient Record'' that constitutes the single, authoritative source of information in the clinic. Participants emphasised that nurses share responsibility for all patients, making the recall of detailed social, medical, and behavioural information challenging without technological support.

Because protected health information was inaccessible, the conversational assistant was evaluated using synthetic but realistic data. Fifty patient records were crafted based on interview insights: the first ten were manually written to mirror authentic documentation styles, while the remainder were generated in \textit{ChatGPT} following the same structure. The dataset mixes structured lists and free-form narratives to represent the variability of clinical notes. No personally identifiable or protected health information was included, ensuring compliance with privacy requirements.

\section{Instruments, Tools, and Software}
Development and experimentation were conducted in Kaggle notebooks equipped with dual Tesla T4 GPUs (15~GB each). The retrieval component employed FAISS for vector search, coupled with the multilingual sentence-transformer \texttt{paraphrase-multilingual-MiniLM-L12-v2} to embed both queries and patient note chunks. Two large language models (LLMs) were explored: the 2~billion parameter ``Gemma\_trained\_for\_Greek'' model and the 8~billion parameter \texttt{ilsp/Llama-Krikri-8B-Instruct}. LangChain provided orchestration utilities, including \texttt{ConversationalRetrievalChain} and \texttt{ConversationBufferMemory}, to inject retrieved context and maintain dialog state. Custom helper functions were introduced to resolve patient names and pronouns when native memory capabilities were insufficient.

\section{Procedures}
The methodological workflow progressed through the following stages:
\begin{enumerate}
    \item \textbf{Requirement elicitation}: synthesising interview findings into functional expectations, such as voice and text interaction, rapid retrieval, and personalisable patient snapshots.
    \item \textbf{Data curation}: transforming interview themes into structured templates, generating synthetic clinical notes, and performing quality checks to ensure coverage of diagnoses, treatments, laboratory schedules, and social histories.
    \item \textbf{Prototype implementation}: configuring the RAG pipeline, designing prompt templates that emphasise factual responses, and integrating memory mechanisms for follow-up queries.
    \item \textbf{Evaluation design}: defining multi-turn scenarios in Greek that reflect daily nurse interactions. The scenarios encompassed happy paths (e.g. medication recall), error handling (unknown patients or missing information), disambiguation cases, free-form summaries, and full-record overviews. Golden responses were authored for each scenario to benchmark faithfulness, relevance, and correctness.
    \item \textbf{Iterative refinement}: analysing system outputs, adjusting chunking strategies to improve patient-specific retrieval, and updating prompt constraints or memory handling when hallucinations appeared.
\end{enumerate}

\section{Justification of Methodological Choices}
Retrieval-augmented generation (RAG) was selected because it combines the generative flexibility of LLMs with traceable, document-grounded responses, mitigating the hallucination risks associated with standalone language models in healthcare \cite{raglewisetal2021}. Agentic orchestration through LangChain enabled tool switching and context management, supporting more natural multi-turn conversations \cite{clinicalentityIE2025}. Synthetic data allowed controlled experimentation without breaching confidentiality, while interviews ensured that mock records reflected the linguistic and informational nuances of Greek psychiatric practice. Evaluating through scenario-based testing provided a transparent link between user needs and observed system behaviour.

\section{Limitations of the Chosen Approach}
The reliance on synthetic records limits external validity until the assistant is trialled with real patient data under clinical supervision. GPU memory ceilings on Kaggle constrained the length of conversations and the size of models that could be deployed, occasionally causing inference crashes after multiple queries. Custom memory solutions depended on accurate pronoun resolution and therefore struggled when multiple patients were referenced in quick succession. Finally, the absence of practicing nurses in the initial evaluation means that perceived usefulness and workflow fit must be validated in future in-clinic studies.

